\section{Impact, use cases, and outlook}
\textbf{Why HCMO matters.} HCM datasets are often distributed as system-specific exports in which measurements are weakly linked—or not linked at all—to the biological subject, enclosure, environmental context, acquisition device, software, and experimental protocol \cite{bains2025,forrest2026,petit2026wellfair}. HCMO provides a shared, machine-actionable semantic model for representing these dependencies consistently. By making the provenance and experimental context of each measurement explicit, HCMO can support dataset discovery, cross-system comparison, integration, and reuse, provided that datasets are mapped to the ontology and exposed with appropriate identifiers and metadata. It thus provides a foundation for transforming otherwise siloed records into interoperable and queryable knowledge. In animal research, where data constitute the durable output of an experiment, preserving their context and interpretability also has ethical relevance: it can reduce avoidable duplication, enable secondary analyses, and increase the scientific value obtained from each animal. HCMO may therefore support the Reduction and Refinement principles of the 3Rs, consistent with the WellFAIR view that data quality and animal welfare are interconnected \cite{petit2026wellfair}.

\textbf{Reasoning and data quality.}
Beyond semantic organisation, HCMO provides a formal basis for automated validation and reasoning. At an immediately applicable level, HCMO and its associated SHACL shapes can identify missing or inconsistent metadata during ingestion---for example, an observation with a numerical value but no unit, or an animal without a valid enclosure assignment during the observation interval. Such checks make data-quality requirements explicit and executable rather than leaving them to manual inspection.

At a second level, HCMO can provide an explicit semantic scaffold for AI-supported analysis. When incorporated into retrieval, query-generation, or agentic workflows, the ontology can specify admissible entity types and relations, as well as the contextual information required to interpret a measurement. The ontology does not, by itself, prevent erroneous inference. Rather, it provides constraints that can be enforced through schema validation, ontology-aware retrieval, and type-constrained query generation, thereby making downstream analyses more controlled, traceable, and interpretable than approaches based on statistical associations alone.

\textbf{Potential uses (outlook).}
HCMO is designed to underpin further tooling that is neither implemented nor evaluated as a contribution of the present work:
\begin{itemize}
    \item \emph{Knowledge-graph question answering.}
    HCMO can provide the schema for an HCM knowledge graph queried using natural language. A user question could be translated into a formal query, checked against the ontology, executed over the graph, and returned with the relevant results and provenance. Implementing this workflow would require reliable natural-language-to-query translation and validation. Related work demonstrates how knowledge graphs and graph embeddings can support structured retrieval, inference, and discovery, including in biomedical applications \cite{sanou2026,bian2025}. The same graph could support symbolic reasoning or graph-learning methods that exploit HCMO's class hierarchy and explicitly typed relations \cite{butz}.

    \item \emph{Assisted authoring.}
    User-facing forms could be generated from the ontology and its SHACL constraints to guide researchers through mandatory and recommended fields. The resulting HCMO-conformant RDF instance data, or ABoxes, could be validated before export or deposition, allowing biologists to produce well-structured data without needing to work directly with RDF, OWL, or SHACL.

    \item \emph{Vendor-data mapping.}
    Mapping components could transform proprietary device exports into HCMO while preserving the original values, identifiers, and provenance. HCMO would thereby provide a common semantic layer for cross-system comparison without requiring vendors or laboratories to adopt a single acquisition format.

\end{itemize}

\textbf{Adoption path.} Development is grounded with a network of preclinical and HCM experts, that initiated with COST-\textbf{TEATIME} action (CA20135), which provides domain feedback and a route to real datasets for validation and uptake. This community anchoring is intended to drive adoption beyond a single laboratory and to keep the model aligned with evolving HCM practices and technologies.
